\documentclass[11pt]{article}
\usepackage[margin=25mm]{geometry}
\usepackage{fontspec}
\setmainfont{Latin Modern Roman}
\usepackage{amsmath,amssymb,mathtools}
\usepackage{unicode-math}
\setmathfont{Latin Modern Math}
\usepackage{xurl,hyperref}
\hypersetup{colorlinks=true,urlcolor=blue}
\setcounter{secnumdepth}{0}
\setlength{\parindent}{0pt}
\setlength{\parskip}{5pt}
\setlength{\emergencystretch}{3em}
\newcommand{\sref}[2]{\hyperlink{source:#1:#2}{[S#2]}}
\newcommand{\cataloguescope}[1]{\paragraph{Recorded target.}#1}
\newcommand{\R}{\mathbb R}
\newcommand{\C}{\mathbb C}
\newcommand{\N}{\mathbb N}
\newcommand{\Z}{\mathbb Z}
\newcommand{\Q}{\mathbb Q}
\newcommand{\F}{\mathbb F}
\newcommand{\eps}{\varepsilon}
\begin{document}
\section{The syzygy-free three-body solution and quantitative crossing certificates}\label{statement}

\subsection{Definitions and mathematical statement}

\paragraph{Shared notation.}$\mathbb N=\{1,2,\ldots\}$; the standard number fields have their usual meanings. All additional conventions are stated locally.


\paragraph{Definitions and notation.}In the Newtonian $n$-body problem, positions $q_i\in\mathbb R^d$ and masses $m_i>0$ obey $m_i\ddot q_i=\nabla_{q_i}U$, where $U(q)=\sum_{i<j}m_im_j/|q_i-q_j|$ and the gravitational constant is normalized to one. Here $T=\tfrac12\sum_i m_i|\dot q_i|^2$, $E=T-U$, and a collision has $q_i=q_j$. The center of mass can be fixed at zero; moment of inertia is $I=\sum_i m_i|q_i|^2$. For the planar three-body problem, a shape is a configuration modulo translation, rotation and positive scaling. Its normalized shape space is a sphere, with collisions marked and collinear configurations on an equator. A syzygy is a collinear configuration; its symbol records which body is in the middle. A central configuration has $\nabla_{q_i}U=-\lambda m_iq_i$ in center-of-mass coordinates for a common scalar $\lambda$. After fixing $I=1$, it is a critical point of $U$ on the collision-free inertia ellipsoid. Congruent configurations and scales are factored out for a finiteness count. A relative equilibrium is a rigidly rotating solution with time-independent shape. A homothetic Lagrange solution has $q_i(t)=r(t)a_i$ with the equilateral central configuration $(a_i)$ and no changing shape. In the planar center-of-mass frame the angular momentum is $C=\sum_i m_iq_i\wedge\dot q_i$. A syzygy-free orbit has no collinear configuration during its maximal collision-free time interval.

\paragraph{Catalogue wording retained for traceability.}
In the three-body problem, are the Lagrange homothety solutions the only solutions with $E<0$, zero angular momentum and no syzygies (eclipses)?
\paragraph{Scope and source qualification.}The maximal interval and treatment of collision endpoints must be made explicit in a proof. The original source question is retained even where the inherited catalogue says solved; it is not newly reclassified here. The catalogue title/status was paired with a machine-generated research attempt; that attempt is not peer-reviewed evidence of settlement. The statement covers the original source questions rather than only the title’s subcase. The inherited catalogue status is solved; this historical statement is retained without independently verifying that settlement.
\subsection{Short English statement}
Determine whether negative-energy, zero-angular-momentum, syzygy-free three-body motion must be Lagrange homothety.
\subsection{Research attempt}
\paragraph{Model, date and investigated scope.}
OpenAI GPT-6 Astra (configured model \texttt{gpt-6-astra}), ultra reasoning; generated 11 October 2026. Verified Lagrange homothety, zero angular momentum, negative-energy turning radius and finite collision endpoints. The source-solved uniqueness/classification is not reconstructed or independently checked.

\paragraph{Sources and literature check.}
Read the supplied complete definition, catalogue wording, scope qualifications and source list for this canonical ID. The original-source citations are inherited as supplied references; no fresh primary-source or novelty search was performed for this attempt. No current open-status claim is inferred from this calculation.

\paragraph{Verifying the exceptional family and its endpoints.}Let \(a_i\) be an equilateral central configuration with center of mass zero, \(I_0=\sum m_i|a_i|^2\), \(U_0=U(a)\), and \(\nabla_iU(a)=-\lambda m_i a_i\). Euler homogeneity gives \(\lambda I_0=U_0\). For \(q_i(t)=r(t)a_i\), \(r>0\), Newton's equations reduce to \(r''=-\lambda/r^2\), and
\[E=\tfrac12I_0\dot r^2-U_0/r,\qquad C=0.
\]
The zero angular momentum follows because each velocity is parallel to its position. The triangle remains equilateral and has nonzero area whenever \(r>0\), so it has no syzygy before a collision endpoint.
\paragraph{Negative energy and finite maximal interval.}For \(E<0\), the maximum radius is \(r_{\max}=U_0/|E|\). The time from radius zero to this turning point is
\[\sqrt{I_0/2}\int_0^{r_{\max}}\frac{dr}{\sqrt{E+U_0/r}},
\]
which is finite: near zero the integrand is \(O(\sqrt r)\), and near the turning point it has an integrable inverse-square-root singularity. Reflection in time gives the collapse half. Thus the syzygy-free example is defined on a finite maximal collision-free interval and ends in total collisions; requiring eternal collision-free motion would exclude the very family named in the source.
\paragraph{Source status and gap.}The catalogue is source-solved, and no fresh verification of its classification proof is claimed. This attempt verifies the exceptional family and the endpoint convention only. To prove exclusivity one must show that every nonhomothetic zero-angular-momentum negative-energy motion has a syzygy on its maximal interval. No shape evolution inequality or global crossing theorem establishing that statement is supplied here; that is the first missing step.

\paragraph{Verification and adversarial audit.}
Checked homogeneity coefficient, radial energy, angular momentum, nonzero area and endpoint integral singularities; preserved source-solved status. This was an LLM adversarial audit, not independent human verification; no theorem prover was used. Prompts A then B applied to this ID: Checked homogeneity coefficient, radial energy, angular momentum, nonzero area and endpoint integral singularities; preserved source-solved status. Scope remains limited: Prove a global syzygy-crossing result for every nonhomothetic zero-angular-momentum negative-energy motion on its maximal collision-free interval.

Coordinator checked the equilateral homothetic equation and finiteness of the radial time integral at collision and turning point; the source classification remains distinct. This additional check was model-assisted, not independent human verification.

\paragraph{Tools, process and contributions.}
Python 3 and local shell used to read supplied files and serialize this manuscript; mathematics derived symbolically in the text, with no numerical experiment or formal proof assistant. Contribution prompt A was followed by prompt B. GPT-6 Astra with ultra reasoning drafted the research and performed its stated self-audit; any additional coordinator contribution is identified in the verification and audit. A GPT-6-based Codex coordinator assembled the manuscript and checked repository integration; its exact variant and reasoning configuration were not exposed. No human mathematical contribution, independent human proof check or formal proof certification is supplied. Complete derivations and the consulted source references are included here; private reasoning transcripts are not published.

\paragraph{Final status and first remaining gap.}
\textbf{UNRESOLVED}. Prove a global syzygy-crossing result for every nonhomothetic zero-angular-momentum negative-energy motion on its maximal collision-free interval.

\section{Completion Estimate}
\noindent\textbf{COMPLETION: 0\%}
The calculation is an elementary benchmark or conditional reduction; it establishes no claimed new portion of the complete source target.

\subsection{Sources}
\begin{itemize}
\item[\hypertarget{source:statement:1}{S1}] ULAM.AI and the UnsolvedMath Contributors, UnsolvedMath, catalogue entry. Catalogue curation under CC BY 4.0; definitions and English exposition added locally.
\url{https://huggingface.co/datasets/ulamai/UnsolvedMath}.
\item[\hypertarget{source:statement:2}{S2}] American Institute of Mathematics, Variational Methods in Celestial Mechanics; notes by Kuo-Chang Chen, version 22 June 2003, source item 16.
\url{https://aimath.org/WWN/varcelest/varcelest.pdf}.
\end{itemize}
\label{end:statement}
\end{document}
